\section{与傅里叶变换有关的其他变换}\label{sec:Other_Transforms}

\subsection{梅林变换}\label{sub:mellin}

\textbf{梅林变换}(Mellin transform)是一种积分变换，定义为：
\begin{definition}[梅林变换]
    设$f\in L^1(\mathbb{R}^+)$，则$f$的梅林变换定义为
    \[\mathcal{M} f(s)=\int_0^{\infty}f(x)x^{s-1}\,dx\]
    其中$s\in\mathbb{C}$。    
\end{definition}

初看起来它与傅里叶变换没什么关系，但如果做变量替换$x=\mathrm{e}^{-t}$，则$dx=-\mathrm{e}^{-t}dt=-x dt$，积分限变为$t\in(-\infty,\infty)$，因此有：
\[\mathcal{M} f(s)=\int_{-\infty}^{\infty}f(\mathrm{e}^{-t})\mathrm{e}^{-s t}\,dt\]
这实际上是$f(\mathrm{e}^{-t})$的双边拉普拉斯变换，见\ref{sub:introduction_of_laplace}，但如果限制$s=\mathrm{i}\omega$，则是$f(\mathrm{e}^{-t})$的角频率形式的傅里叶变换：
\[\mathcal{M} f(\mathrm{i}\omega)=\int_{-\infty}^{\infty}f(\mathrm{e}^{-t})\mathrm{e}^{-\mathrm{i}\omega t}\,dt\]
可以通过傅里叶变换的性质研究梅林变换的性质。显然，梅林变换是线性的。

首先，利用傅里叶反演公式（见\ref{sub:introduction_of_fourier}或\ref{sub:fourier_inversion}），可以得到梅林变换的反变换：
\begin{proposition}[梅林反变换]
    设$f\in L^1(\mathbb{R}^+)$，则有
    \[f(x)=\frac{1}{2\pi \mathrm{i}}\int_{-\infty}^{\infty}\mathcal{M} f(\mathrm{i}\omega)x^{-\mathrm{i}\omega}\,d\omega\]
\end{proposition}
\begin{proof}
    由傅里叶反演公式，有
    \[f(\mathrm{e}^{-t})=\mathcal{F} ^{-1}[\mathcal{M} f(\mathrm{i}\omega)](t)=\frac{1}{2\pi}\int_{-\infty}^{\infty}\mathcal{M} f(\mathrm{i}\omega)\mathrm{e}^{\mathrm{i}\omega t}\,d\omega\]
    令$t=-\ln x$，则有
    \[f(x)=\frac{1}{2\pi}\int_{-\infty}^{\infty}\mathcal{M} f(\mathrm{i}\omega)\mathrm{e}^{\mathrm{i}\omega (-\ln x)}\,d\omega=\frac{1}{2\pi}\int_{-\infty}^{\infty}\mathcal{M} f(\mathrm{i}\omega)x^{-\mathrm{i}\omega}\,d\omega\]
\end{proof}

在学习\ref{sub:laplace_inversion}后，也可以用类似的方法，将梅林反变换的求解转化为留数的计算，不过我们不打算深入讨论这个话题。

梅林变换可以看作是傅里叶变换在对数坐标系下的对应。从这个视角来看，梅林变换在时域发生变化时，频域的变化方式会与傅里叶变换有所不同。具体来说，我们给出以下命题：
\begin{proposition}[梅林变换的尺度变换性质]
    设$f\in L^1(\mathbb{R}^+)$，则对于任意$a>0$，都有
    \[\mathcal{M} [f(ax)](s)=a^{-s}\mathcal{M} f(s)\]
\end{proposition}
\begin{proof}
    由梅林变换的定义，有
    \begin{align*}
        \mathcal{M} [f(ax)](s)&=\int_0^{\infty}f(ax)x^{s-1}\,dx\\
        &=\int_0^{\infty}f(u)a^{-1}(u/a)^{s-1}\,du\quad (u=ax,dx=du/a)\\
        &=a^{-s}\int_0^{\infty}f(u)u^{s-1}\,du\\
        &=a^{-s}\mathcal{M} f(s)
    \end{align*}
\end{proof}

这表明时域上的尺度变换对应频域上的指数变换，这与傅里叶变换中时域的尺度变换对应频域的反比例变换不同。此外，梅林变换没有时移性质，不过我们还是可以提出变换域上的平移性质：

\begin{proposition}[梅林变换的乘幂性质]
    设$f\in L^1(\mathbb{R}^+)$，则对于任意$c\in\mathbb{C}$，都有
    \[\mathcal{M} [x^c f(x)](s)=\mathcal{M} f(s+c)\]
\end{proposition}
\begin{proof}
    由梅林变换的定义，有
    \begin{align*}
        \mathcal{M} [x^c f(x)](s)&=\int_0^{\infty}x^c f(x)x^{s-1}\,dx\\
        &=\int_0^{\infty}f(x)x^{s+c-1}\,dx\\
        &=\mathcal{M} f(s+c)
    \end{align*}
\end{proof}

对于傅里叶变换，频域或角频域上的平移对应时域上的调制，即乘以某个复指数信号；而梅林变换中，频域上的平移对应时域上的乘幂。

此外，梅林变换还有卷积性质，不过这里的卷积不是通常意义下的卷积，而是\textbf{梅林卷积}，定义为：
\begin{definition}[梅林卷积]
    设$f,g\in L^1(\mathbb{R}^+)$，则它们的\textbf{梅林卷积}(Mellin convolution)定义为
    \[(f\star g)(x)=\int_0^{\infty}f\left(\frac{x}{t}\right)g(t)\frac{dt}{t}\]
\end{definition}

符号$\star$与\ref{sub:correlation}中的互相关符号相同，但我们很少同时用到它们，必要时将指出它的具体含义。

%我们还将在\ref{sub:z_convolution}看到类似的卷积方式，并且梅林变换的定义与性质与z变换有一定的相似之处，不过，z变换是用来处理离散信号的，而梅林卷积处理的是连续信号。它们的相似性或许可以解释为，它们都可以看作是特定情况下的双边拉普拉斯变换。

\begin{proposition}[梅林卷积的变换性质]
    设$f,g\in L^1(\mathbb{R}^+)$，则有
    \[\mathcal{M} [f\star g](s)=\mathcal{M} f(s)\cdot \mathcal{M} g(s)\]
\end{proposition}
\begin{proof}
    由梅林卷积的定义，有
    \begin{align*}
        \mathcal{M} [f\star g](s)&=\int_0^{\infty}(f\star g)(x)x^{s-1}\,dx\\
        &=\int_0^{\infty}\int_0^{\infty}f\left(\frac{x}{t}\right)g(t)\frac{dt}{t}x^{s-1}\,dx\\
        &=\int_0^{\infty}g(t)\frac{1}{t}\int_0^{\infty}f\left(\frac{x}{t}\right)x^{s-1}\,dx\,dt\\
        &=\int_0^{\infty}g(t)\frac{1}{t}\int_0^{\infty}f(u)(ut)^{s-1}t\,du\,dt\quad (u=x/t,dx=t\,du)\\
        &=\int_0^{\infty}g(t)t^{s-1}\,dt\int_0^{\infty}f(u)u^{s-1}\,du\\
        &=\mathcal{M} f(s)\cdot \mathcal{M} g(s)
    \end{align*}
\end{proof}

\subsection{希尔伯特变换}\label{sub:hilbert}

\textbf{希尔伯特变换}(Hilbert transform)既可以看作积分变换，又可以看作滤波器，因为它是函数与$\dfrac{1}{\pi t}$的卷积（见\ref{sub:cascade}）。希尔伯特变换在信号处理领域有着重要的应用，例如构造解析信号，将在本小节最后介绍。
\begin{definition}[希尔伯特变换]
    函数$f:\mathbb{R}\to\mathbb{C}$的\textbf{希尔伯特变换}定义为
    \[\mathcal{H} f(x)=\frac{1}{\pi}\text{p.v.}\int_{-\infty}^{\infty}\frac{f(y)}{x-y}\,dy=\left(f*\frac{1}{\pi x}\right)(x)\]
    其中p.v.表示柯西主值意义下的积分\cite{hilbert_transform}，即：
    \[\text{p.v.}\int_{-\infty}^{\infty}\frac{f(y)}{x-y}\,dy=\lim_{\epsilon\to 0^+}\left(\int_{-\infty}^{x-\epsilon}\frac{f(y)}{x-y}\,dy+\int_{x+\epsilon}^{\infty}\frac{f(y)}{x-y}\,dy\right)\]
    称$f$与$\mathcal{H} f$为\textbf{希尔伯特变换对}，记为$f\overset{\mathcal{H} }{\longleftrightarrow} \mathcal{H} f$。
\end{definition}
\begin{remark}
    一些文献中\cite{brad_osgood_ee261}，希尔伯特变换为$-\dfrac{1}{\pi t}$与$f(t)$的卷积，我们很快将看到这就是本书中的希尔伯特逆变换，它们没有本质区别；一些文献中\cite{signal_systems}，直接将希尔伯特变换记为$\hat{f}$，这样做有一定的好处，例如书写简便，但它与傅里叶变换的记号冲突，因此本书中不采用这种记号。
\end{remark}
\begin{remark}
    我们没有对$f$所属的函数空间做出限制。实际上，希尔伯特变换可以定义在多种函数空间上，经典的希尔伯特变换建立在$L^1(\mathbb{R})$上，并且要求函数$f$是赫德尔连续的（我们没有介绍这个概念），希尔伯特变换还可以定义在$L^p(\mathbb{R}),1<p<\infty$上，或者更广义的分布空间上。专著\cite{hilbert_transform}详细讨论了$f$的范围。
\end{remark}

在\ref{sub:fourier_of_signal}中，我们求出了函数$\dfrac{1}{x}$的傅里叶变换，并在建立了分布的理论后能够严格证明：
\begin{lralign}
    \mathcal{F} \left[\frac{1}{\pi x}\right](\xi)= -\mathrm{i}\,\text{sgn}(\xi)
    \switch
    \mathcal{F} \left[\frac{1}{\pi x}\right](\omega)= -\mathrm{i}\,\text{sgn}(\omega)
\end{lralign}

借助这个结论，我们可以通过傅里叶变换来研究希尔伯特变换的性质。简洁起见，本小节主要使用频率形式的傅里叶变换。

由卷积定理（见\ref{sub:properties_of_convolution}），我们可以得到希尔伯特变换的傅里叶变换表达式：
    \[\mathcal{F} [\mathcal{H} f](\xi)=-\mathrm{i}\,\text{sgn}(\xi)\mathcal{F} f(\xi)\]
可以看到，希尔伯特变换在频域上对应乘以$-\mathrm{i}\,\text{sgn}(\xi)$，这表明希尔伯特变换是一个全通滤波器，它将正频率成分的相位延迟$\pi/2$，将负频率成分的相位超前$\pi/2$，而不改变任意频率成分的幅度。

如果连续取两次希尔伯特变换，在频域上就对应乘以$(-\mathrm{i}\,\text{sgn}(\xi))^2=-1$，也就是：
\[\mathcal{H} \mathcal{H} f=-f\]
因此，希尔伯特变换的逆变换为：
\begin{definition}[希尔伯特逆变换]
    设$f\in L^1(\mathbb{R})$，则$f$的\textbf{希尔伯特逆变换}定义为
    \[\mathcal{H} ^{-1} f(x)=-\mathcal{H} f(x)=-\frac{1}{\pi}\text{p.v.}\int_{-\infty}^{\infty}\frac{f(y)}{x-y}\,dy=-\left(f*\frac{1}{\pi x}\right)(x)\]
\end{definition}
至此可以看出，希尔伯特变换对是一一的，因此记号$f\overset{\mathcal{H} }{\longleftrightarrow} \mathcal{H} f$是合理的。

如果$f\in L^2(\mathbb{R})$，即$f(t)$是一个能量有限信号，则$\mathcal{H} f\in L^2(\mathbb{R})$，并且$\|\mathcal{H} f\|_2=\|f\|_2$，进一步，$L^2(\mathbb{R})$上的\textbf{希尔伯特变换是保内积的}：
\begin{theorem}
    设$f,g\in L^2(\mathbb{R})$，则有
    \[(\mathcal{H} f,\mathcal{H} g)=(f,g)\]
    特别地，$\|\mathcal{H} f\|_2=\|f\|_2$，这称之为希尔伯特变换的\textbf{等能量性}。
\end{theorem}
\begin{proof}
    由希尔伯特变换的傅里叶变换表达式，有
    \begin{align*}
        (\mathcal{H} f,\mathcal{H} g)&=\int_{-\infty}^{\infty}\mathcal{H} f(t)\overline{\mathcal{H} g(t)}\,dt\\
        &=\int_{-\infty}^{\infty}\mathcal{F} [\mathcal{H} f](\xi)\overline{\mathcal{F} [\mathcal{H} g](\xi)}\,d\xi\\
        &=\int_{-\infty}^{\infty}(-\mathrm{i}\,\text{sgn}(\xi)\mathcal{F} f(\xi))\overline{-\mathrm{i}\,\text{sgn}(\xi)\mathcal{F} g(\xi)}\,d\xi\\
        &=\int_{-\infty}^{\infty}\mathcal{F} f(\xi)\overline{\mathcal{F} g(\xi)}\,d\xi\\
        &=(f,g)
    \end{align*}
\end{proof}

希尔伯特变换还具有\textbf{正交性}：
\begin{theorem}
    设$f\in L^2(\mathbb{R})$，则有
    \[(f,\mathcal{H} f)=0\]
\end{theorem}
由于这个原因，有时将$\mathcal{H} f$称为$f$的\textbf{正交伴随}(orthogonal complement)。
\begin{proof}
    由希尔伯特变换的傅里叶变换表达式，有
    \begin{align*}
        (f,\mathcal{H} f)&=\int_{-\infty}^{\infty}f(t)\overline{\mathcal{H} f(t)}\,dt\\
        &=\int_{-\infty}^{\infty}\mathcal{F} f(\xi)\overline{\mathcal{F} [\mathcal{H} f](\xi)}\,d\xi\\
        &=\int_{-\infty}^{\infty}\mathcal{F} f(\xi)\overline{-\mathrm{i}\,\text{sgn}(\xi)\mathcal{F} f(\xi)}\,d\xi\\
        &=-\mathrm{i}\int_{-\infty}^{\infty}\text{sgn}(\xi)|\mathcal{F} f(\xi)|^2\,d\xi\\
        &=0
    \end{align*}
\end{proof}

此外，希尔伯特变换具有奇偶对称性：
\begin{theorem}
    设$f\in L^2(\mathbb{R})$，则有
    \begin{itemize}
        \item 如果$f$是偶函数，则$\mathcal{H} f$是奇函数；
        \item 如果$f$是奇函数，则$\mathcal{H} f$是偶函数。
    \end{itemize}
\end{theorem}
\begin{proof}
    设$f$是偶函数，则$\mathcal{F} f$也是偶函数，因此$\mathcal{F} [\mathcal{H} f](\xi)=-\mathrm{i}\,\text{sgn}(\xi)\mathcal{F} f(\xi)$是奇函数，因此$\mathcal{H} f$是奇函数。设$f$是奇函数，则$\mathcal{F} f$也是奇函数，因此$\mathcal{F} [\mathcal{H} f](\xi)=-\mathrm{i}\,\text{sgn}(\xi)\mathcal{F} f(\xi)$是偶函数，因此$\mathcal{H} f$是偶函数。
\end{proof}

\begin{example}[三角函数的希尔伯特变换]
    三角函数的希尔伯特变换可以通过它们的傅里叶变换来求解。考虑$\cos(2\pi f_0 x)=\cos(\omega_0 x)$，我们知道：
    \begin{lralign}
        \mathcal{F} [\cos(2\pi f_0 t)](\xi)&=\frac{1}{2}(\delta_{f_0}+\delta_{-f_0})
        \switch
        \mathcal{F} [\cos(\omega_0 t)](\xi)&=\pi(\delta_{\omega_0}+\delta_{-\omega_0})
    \end{lralign}
    因此余弦函数的希尔伯特变换的傅里叶变换为：
    \begin{lralign}
        &\mathcal{F} [\mathcal{H} \cos(2\pi f_0 t)](\xi)\\
        =&-\mathrm{i}\,\text{sgn}(\xi)\cdot \frac{1}{2}(\delta_{f_0}+\delta_{-f_0})\\
        =&-\frac{\mathrm{i}}{2}(\delta_{f_0}-\delta_{-f_0})
        \switch
        &\mathcal{F} [\mathcal{H} \cos(\omega_0 t)](\xi)\\
        =&-\mathrm{i}\,\text{sgn}(\xi)\cdot \pi(\delta_{\omega_0}+\delta_{-\omega_0})\\
        =&-\mathrm{i}\pi(\delta_{\omega_0}-\delta_{-\omega_0})
    \end{lralign}
    这正是$\sin(2\pi f_0 t)=\sin(\omega_0 t)$的傅里叶变换，因此我们得到：
    \[\mathcal{H} \cos(\omega_0 t)=\sin(\omega_0 t)\]
    同理可得：
    \[\mathcal{H} \sin(\omega_0 t)=-\cos(\omega_0 t)\]
\end{example}

\begin{example}[直接计算希尔伯特变换]
    对于多数函数而言，计算希尔伯特变换时，直接使用定义进行计算比较困难。这里给出一个简单的函数，它能够使用定义计算希尔伯特变换。考虑脉冲宽度为$T$的矩形函数$f(x)=\chi[-T/2,T/2](x)$，则有：
    \begin{align*}
        \mathcal{H} f(x)&=\frac{1}{\pi}\text{p.v.}\int_{-T/2}^{T/2}\frac{1}{x-y}\,dy\\
        &=\frac{1}{\pi}\lim_{\epsilon\to 0^+}\left(\int_{-T/2}^{x-\epsilon}\frac{1}{x-y}\,dy+\int_{x+\epsilon}^{T/2}\frac{1}{x-y}\,dy\right)\\
        &=\frac{1}{\pi}\lim_{\epsilon\to 0^+}\left(\ln|x-y|\Big|_{y=-T/2}^{y=x-\epsilon}+\ln|x-y|\Big|_{y=x+\epsilon}^{y=T/2}\right)\\
        &=\frac{1}{\pi}\lim_{\epsilon\to 0^+}\left(\ln \epsilon -\ln|x+T/2|+\ln|T/2 -x|-\ln \epsilon \right)\\
        &=\frac{1}{\pi}\ln\left|\frac{T/2 -x}{x+T/2}\right|
    \end{align*}
\end{example}

\begin{example}[三角函数的希尔伯特变换，直接计算]
    如果使用狄利克雷积分的结果：
    \[\int_{-\infty}^{\infty}\frac{\sin ax}{x}\,dx=\pi,a>0\]
    也可以直接计算三角函数的希尔伯特变换，这次以正弦函数为例：
    \begin{align*}
        \mathcal{H} \sin(\omega_0 t)&=\frac{1}{\pi}\text{p.v.}\int_{-\infty}^{\infty}\frac{\sin(\omega_0 y)}{t-y}\,dy\\
        &=\frac{1}{\pi}\lim_{\epsilon\to 0^+}\left(\int_{-\infty}^{t-\epsilon}\frac{\sin(\omega_0 y)}{t-y}\,dy+\int_{t+\epsilon}^{\infty}\frac{\sin(\omega_0 y)}{t-y}\,dy\right)\\
        &=\frac{1}{\pi}\lim_{\epsilon\to 0^+}\left(\int_{\epsilon}^{\infty}\frac{\sin(\omega_0 (t-u))}{u}\,du+\int_{-\infty}^{-\epsilon}\frac{\sin(\omega_0 (t-u))}{u}\,du\right)& (u=t-y)\\
        &=\frac{1}{\pi}\lim_{\epsilon\to 0^+}\left(\int_{\epsilon}^{\infty}\frac{\sin(\omega_0 t)\cos(\omega_0 u)-\cos(\omega_0 t)\sin(\omega_0 u)}{u}\,du\right.\\
        &\quad \left.+\int_{-\infty}^{-\epsilon}\frac{\sin(\omega_0 t)\cos(\omega_0 u)-\cos(\omega_0 t)\sin(\omega_0 u)}{u}\,du\right)\\
        &=\frac{1}{\pi}\lim_{\epsilon\to 0^+}\left(\sin(\omega_0 t)\int_{\epsilon}^{\infty}\frac{\cos(\omega_0 u)}{u}\,du+\sin(\omega_0 t)\int_{-\infty}^{-\epsilon}\frac{\cos(\omega_0 u)}{u}\,du\right.\\
        &\quad \left.-\cos(\omega_0 t)\int_{\epsilon}^{\infty}\frac{\sin(\omega_0 u)}{u}\,du-\cos(\omega_0 t)\int_{-\infty}^{-\epsilon}\frac{\sin(\omega_0 u)}{u}\,du\right)\\
        &=\frac{1}{\pi}\lim_{\epsilon\to 0^+}\left(-\cos(\omega_0 t)\int_{\epsilon}^{\infty}\frac{\sin(\omega_0 u)}{u}\,du+\cos(\omega_0 t)\int_{\epsilon}^{\infty}\frac{\sin(\omega_0 u)}{u}\,du\right)\\
        &=-\cos(\omega_0 t)
    \end{align*}
    倒数第三行到倒数第二行运用了函数的奇偶性。
\end{example}

希尔伯特变换在信号分析中有着广泛的应用，下面给出由它得到的一些重要结果。
\begin{definition}[解析信号]
    函数$f\in L^1(\mathbb{R})$称为\textbf{解析信号}(analytic signal)，如果它仅含正频率成分，即
    \[\text{supp}\mathcal{F} f\subseteq[0,\infty)\]
\end{definition}

根据傅里叶变换的对偶性（见\ref{sub:fourier_operation_properties}），实信号的傅里叶变换满足$\mathcal{F} f^-=\overline{\mathcal{F} f}$，即幅度谱为偶函数而相位谱为奇函数，因此实信号不可能是解析信号。但如果考虑$2u\cdot\mathcal{F} f$，其中$u$为单位阶跃函数，则该信号仅含正频率成分（我们马上将看到为什么是两倍的$u\mathcal{F} f$）。

对$2u\cdot\mathcal{F} f$做傅里叶逆变换，并即相应的时域函数为$\mathcal{Z}_f(t)$，则有：
\begin{lralign}
    \mathcal{Z}_f(t)&=\mathcal{F} ^{-1}[2u(\xi)\cdot\mathcal{F} f(\xi)](t)\\
    &=\mathcal{F} ^{-1}[(1+\text{sgn}(\xi))\mathcal{F} f(\xi)](t)\\
    &=\mathcal{F} \mathcal{F} ^{-1}f(t)+\mathcal{F} ^{-1}[\text{sgn}(\xi)\mathcal{F} f(\xi)](t)\\
    &=f(t)+\mathrm{i}\mathcal{H} f(t)
    \switch
    \mathcal{Z}_f(t)&=\mathcal{F} ^{-1}[2u(\omega)\cdot\mathcal{F} f(\omega)](t)\\
    &=\mathcal{F} ^{-1}[(1+\text{sgn}(\omega))\mathcal{F} f(\omega)](t)\\
    &=\mathcal{F} \mathcal{F} ^{-1}f(t)+\mathcal{F} ^{-1}[\text{sgn}(\omega)\mathcal{F} f(\omega)](t)\\
    &=f(t)+\mathrm{i}\mathcal{H} f(t)
\end{lralign}

因此，在实信号的虚部中加入它的希尔伯特变换（显然，实信号的希尔伯特变换仍是实信号），就得到了一个解析信号。需要注意的是，以上解析信号的构造过程并未限制$f$是实信号。我们定义：

\begin{definition}
    设$f\in L^1(\mathbb{R})$，称$\mathcal{Z}_f(t)=f(t)+\mathrm{i}\mathcal{H} f(t)$为\textbf{与}$f$\textbf{相关的解析信号}。
    从构造过程可以看出，时域上，$f(x)=\text{Re}\{\mathcal{Z}_f(x)\}$；频域上，$\mathcal{Z} _f$是由$f$截去复频率成分，而正频率成分加倍得到的。
\end{definition}

与解析信号相关的解析信号并不是它本身，因为只有0的希尔伯特变换是0。

以上构造过程告诉我们，如果一个线性时不变系统能够的响应是激励信号的希尔伯特变换，就能够用它获得与激励信号相关的解析信号。但是，希尔伯特变换的单位冲激响应为$\dfrac{1}{\pi t}$，它不是因果信号，因此希尔伯特变换系统不是因果系统。

最后，我们来讨论因果的线性时不变系统的性质。设$h(t)$为因果线性时不变系统的单位冲激响应，则$h(t)u(t)=h(t)$。考虑系统的频率响应$H(\omega)=\mathcal{F} h(\omega)$，则有：
\begin{align*}
    H(\omega)&=\mathcal{F} [h(t)u(t)](\omega)\\
    &=\frac{1}{2\pi}H(\omega)*\left(\pi\delta(\omega)+\frac{1}{\mathrm{i}\omega}\right)\\
    &=\frac{1}{2}H(\omega)+\frac{1}{2\pi \mathrm{i}\omega}*H(\omega)
\end{align*}

考察上式的实部和虚部，记$H(\omega)=R(\omega)+iI(\omega)$，则有：
\begin{align*}
    R(\omega)&=\frac{1}{2}R(\omega)-\frac{1}{2\pi \omega}*I(\omega)\\
    I(\omega)&=\frac{1}{2}I(\omega)+\frac{1}{2\pi \omega}*R(\omega)
\end{align*}
整理可得：
\begin{align*}
    &R(\omega)=-\frac{1}{\pi \omega}*I(\omega)=\mathcal{H} ^{-1}I(\omega)\\
    &I(\omega)=\frac{1}{\pi \omega}*R(\omega)=\mathcal{H} R(\omega)
\end{align*}
也就是以下定理：
\begin{theorem}[希尔伯特变换关系]
    设$h(t)$为因果线性时不变系统的单位冲激响应，$H(\omega)=R(\omega)+iI(\omega)$为系统的频率响应，则$R(\omega)$与$I(\omega)$构成希尔伯特变换对，即：
    \begin{align*}
    &R(\omega)=-\frac{1}{\pi \omega}*I(\omega)=\mathcal{H} ^{-1}I(\omega)\\
    &I(\omega)=\frac{1}{\pi \omega}*R(\omega)=\mathcal{H} R(\omega)
\end{align*}
将这种关系称为\textbf{希尔伯特变换关系}(Hilbert transform relation)。
\end{theorem}

\subsection{拉东变换}\label{sub:radon}

拉东变换在医疗成像领域有着重要的应用，它使得我们可以通过x射线对人体进行扫描，进而重建人体内部的结构图像。本小节将介绍x射线成像的基本原理，读者将看到，一些资料中将这种技术描述为“x光透视”是不准确的。

考虑一束光强为$I_0$的光线穿过浑浊的水，则光线强度会随穿过水的距离增加而减小，如果浑浊程度是均匀的，我们可以用$I=I_0 \mathrm{e}^{-\mu x}$来描述光线强度的变化，其中$x$为光线穿过水的距离，$\mu$描述了水的浑浊程度。如果水的浑浊程度不均匀，则在光线所在线段$L$上，$\mu$是空间位置的函数$\mu(x)$，出射光线的强度为：
\[I=I_0 \exp\left(-\int_{L}\mu(x)\,dx\right)\]

类似地，我们考虑一个水平移动的x射线源穿过人体，对人体进行扫描，而另一侧为探测器阵列，用于测量穿过人体后x射线的强度变化。这样，每一次扫描，我们都将获得若干直线上的x射线的总衰减量，可以想象，在适当的数据处理下，我们可以重建人体内部的衰减系数分布$\mu(x,y)$，进而得到人体内部的二维截面的结构图像。这就是\textbf{拉东变换}(Radon transform)做的工作。

拉东变换是一种广义的积分变换。积分变换通常将函数与某个权函数（例如，傅里叶变换时，权函数为$\mathrm{e}^{-2\pi \mathrm{i}\xi t}$或$\mathrm{e}^{-\mathrm{i}\omega t}$）相乘再积分，而拉东变换的则将函数与“权分布”相乘（这并不是一个标准的术语）。在给出它的定义之前，我们先来约定本小节中的符号。

\begin{definition}[二位平面上的直线系]
    二维平面上的直线$L:ax+by+c=0$可以由参数组$(\rho,\phi)$唯一确定，其中$\rho\geq 0$为原点到直线的距离，$\phi\in[0,\pi)$为$L$与x轴正方向的夹角。
    我们将$L$记为$L(\rho,\phi)$。
\end{definition}

在这样的符号下，$L(\rho,\phi)$具有单位法向量$\hat{\mathbf{n}}=(\cos\phi,\sin\phi)$，并且$$\forall (x,y)\in L(\rho,\phi),\rho=\mathbf{x}\cdot\hat{\mathbf{n}}=x\cos\phi+y\sin\phi$$
因此$L(\rho,\phi)$的方程为：
\[L(\rho,\phi): x\cos\phi+y\sin\phi-\rho=0\]

本小节中，我们总是将这个方程记为$l(x,y)=0$，其中函数$l(x,y)=x\cos\phi+y\sin\phi-\rho$。

\begin{definition}[线脉冲]
    设$L(\rho,\phi)$为二维平面上的一条直线，则它的\textbf{线脉冲}(line delta function)定义为
    \[\delta_{L(\rho,\phi)}=\delta(l(x,y))=\delta(x\cos\phi+y\sin\phi-\rho)\]
    其中$\delta$为一维狄拉克函数。尽管它只是一个形式上的函数，但借助它来理解拉东变换是方便的。
\end{definition}

以上我们给出了线脉冲在形式上的定义，作为分布而言，线脉冲的作用为：
\[\left\langle \delta_{L(\rho,\phi)},f\right\rangle=\int_{\mathbb{R}^2}\delta(l(x,y))f(x,y)\,dx\,dy=\int_{L(\rho,\phi)}f(x,y)\,ds\]
其中$ds$为弧长微元。也就是说，线脉冲作为泛函作用于二元函数$f(x,y)$，相当于$f$在直线$L(\rho,\phi)$上的线积分。实际上，它是以下定理的推论，不过我们不在这里进行证明。
\begin{claim}
    设有平面曲线$\gamma:F(x,y)=0$，$F$连续可微，且$|\nabla F|\neq 0$，其中
    \[\nabla F=\text{grad}F=\begin{bmatrix}
        \frac{\partial F}{\partial x}\\
        \frac{\partial F}{\partial y}        
    \end{bmatrix}\]
    为$F$的梯度，$|\cdot|$为向量值函数$\nabla F$的模长，则有：
    \begin{align*}
        \left\langle \delta_{\gamma},f\right\rangle=\int_{\mathbb{R}^2}\delta(F(x,y))f(x,y)\,dx\,dy=\int_{\gamma}\frac{f(x,y)}{|\nabla F|}\,ds
    \end{align*}
\end{claim}

类似的，线脉冲可以对高维空间中的\textbf{超平面}(hyperplane)定义。
\begin{definition}
$\mathbb{R}^n$中的超平面为：
\[P=\left\{\mathbf{x}\in\mathbb{R}^n\left|a_1 x_1+a_2 x_2+\cdots+a_n x_n=0,a_j(1\leq j\leq n)\text{为常数}\right.\right\}\]
在三维空间中，它就是我们所熟知的平面，而在n维空间中为n-1维流形，可以用参数对$(\rho,\hat{\mathbf{n}})$表示，其中$\rho$为原点到超平面的距离，$\hat{\mathbf{n}}$为单位法向量。我们又将超平面$P$记为$P(\rho,\hat{\mathbf{n}})$。尽管这样使用了n个参量，但$\hat{\mathbf{n}}$的参量不是独立的。
\end{definition}

现在，终于可以给出拉东变换的定义。
\begin{definition}[拉东变换]
    给定平面上的二元函数$f(x,y)$，它的拉东变换定义为：
    \begin{align*}
    \mathcal{R} f(\rho,\phi)=\left\langle \delta_{L(\rho,\phi)},f\right\rangle=\int_{\mathbb{R}^2}\delta(l(x,y))f(x,y)\,dx\,dy=\int_{L(\rho,\phi)}f(x,y)\,ds
    \end{align*}
\end{definition}
\begin{remark}
    我们也没有对$f$的范围进行讨论。一般而言，我们需要要求$f$分段连续，并且是紧支的或至少在无穷远处快速衰减至0的（例如施瓦兹函数），这样才能保证线积分存在。这个函数空间的精确描述由Cormack定理给出，我们不对它进行讨论，仅仅指出：这个函数空间构成线性空间。
\end{remark}
\begin{remark}
    沿用前面对超平面的记号，对于n元函数$f(\mathbf{x})$及超平面$P(\rho,\hat{\mathbf{n}})$，拉东变换定义为
    \begin{align*}
    \mathcal{R} f(\rho,\hat{\mathbf{n}})=\int_{\mathbb{R}^n}f(\mathbf{x})\delta(\rho-\mathbf{x}\cdot\hat{\mathbf{n}})\,d\mathbf{x}=\int_{P}f(\mathbf{x})\,d\Omega
    \end{align*}
    其中$d\Omega$是$n-1$维流形上的体积元。高维的拉东变换不是本小节的主要目标。
\end{remark}

\begin{example}[圆形特征函数的拉东变换]
    设有一以原点为中心，半径为$r$的圆形区域$$D=\{(x,y)|x^2+y^2\leq r^2\}$$
    考虑它的特征函数$\chi_D(x,y)$，则它的拉东变换为：
    \begin{align*}
        \mathcal{R} \chi_D(\rho,\phi)&=\int_{\mathbb{R}^2}\chi_D(x,y)\delta(x\cos\phi+y\sin\phi-\rho)\,dx\,dy\\
        &=\int_{L(\rho,\phi)}\chi_D(x,y)\,ds
    \end{align*}
    当$|\rho|>r$时，直线$L(\rho,\phi)$与圆形区域$D$没有交点，因此$\mathcal{R} \chi_D(\rho,\phi)=0$。当$|\rho|\leq r$时，直线$L(\rho,\phi)$与圆形区域$D$有两个交点，记为$A,B$，则：
    \begin{align*}
        \mathcal{R} \chi_D(\rho,\phi)&=\int_{L(\rho,\phi)}\chi_D(x,y)\,ds=|AB|=2\sqrt{r^2-\rho^2}
    \end{align*}
    综上所述，圆形区域的特征函数的拉东变换为：
    \[\mathcal{R} \chi_D(\rho,\phi)=\begin{cases}
        2\sqrt{r^2-\rho^2},&|\rho|\leq r\\
        0,&|\rho|>r
    \end{cases}\]
\end{example}
\begin{example}[高斯函数的拉东变换]
    二维高斯函数$\mathrm{e}^{-\pi(x^2+y^2)}=\mathrm{e}^{-\pi|\mathbf{x}|^2}$的拉东变换为：
    \begin{align*}
        \mathcal{R} [\mathrm{e}^{-\pi|\mathbf{x}|^2}](\rho,\phi)&=\int_{\mathbb{R}^2}\mathrm{e}^{-\pi(x^2+y^2)}\delta(x\cos\phi+y\sin\phi-\rho)\,dx\,dy\\
        &=\int_{L(\rho,\phi)}\mathrm{e}^{-\pi(x^2+y^2)}\,ds
    \end{align*}
    做正交变量替换：
    \[\begin{bmatrix}u\\v\end{bmatrix}=\begin{bmatrix}
        \cos\phi & \sin\phi\\
        -\sin\phi & \cos\phi
    \end{bmatrix}\begin{bmatrix}x\\y\end{bmatrix}\]
    则有$ds=dv$，并且在直线$L(\rho,\phi)$上，$u=\rho$，因此：
    \begin{align*}
        \mathcal{R} [\mathrm{e}^{-\pi|\mathbf{x}|^2}](\rho,\phi)&=\int_{-\infty}^{\infty}\mathrm{e}^{-\pi(\rho^2+v^2)}\,dv\\
        &=\mathrm{e}^{-\pi \rho^2}\int_{-\infty}^{\infty}\mathrm{e}^{-\pi v^2}\,dv\\
        &=\mathrm{e}^{-\pi \rho^2}
    \end{align*}
    高斯函数的拉东变换仍是高斯函数，令人惊讶。
\end{example}

下面，我们首先给出一些二维拉东变换的基本性质，再使用傅里叶变换来研究拉东变换。

\begin{proposition}[线性性]
    拉东变换是线性的，即任给$f,g:\mathbb{R}^2\to\mathbb{C}$及常数$a,b\in\mathbb{C}$，有
    \[\mathcal{R}[af+bg](\rho,\phi)=a\mathcal{R} f(\rho,\phi)+b\mathcal{R} g(\rho,\phi)\]
\end{proposition}
这个性质是显然的。

\begin{proposition}[平移性质]
    \[\mathcal{R} [\tau_{\mathbf{b}}f](\rho,\hat{\mathbf{n}})=\mathcal{R} f(\rho-\hat{\mathbf{n}}\cdot\mathbf{b},\hat{\mathbf{n}})\]
\end{proposition}
$\rho-\hat{\mathbf{n}}\cdot\mathbf{b}<0$的情况，将在下一性质指出处理方法。
\begin{proof}
    \begin{align*}
        \mathcal{R} [\tau_{\mathbf{b}}f](\rho,\hat{\mathbf{n}})&=\int_{\mathbb{R}^2}f(\mathbf{x}-\mathbf{b})\delta(\rho-\mathbf{x}\cdot\hat{\mathbf{n}})\,d\mathbf{x}\\
        &=\int_{\mathbb{R}^2}f(\mathbf{y})\delta(\rho-(\mathbf{y}+\mathbf{b})\cdot\hat{\mathbf{n}})\,d\mathbf{y}&(\mathbf{y}=\mathbf{x}-\mathbf{b})\\
%        &=\int_{\mathbb{R}^2}f(\mathbf{y})\delta((\rho-\mathbf{b}\cdot\hat{\mathbf{n}})-\mathbf{y}\cdot\hat{\mathbf{n}})\,d\mathbf{y}\\
        &=\mathcal{R} f(\rho-\hat{\mathbf{n}}\cdot\mathbf{b},\hat{\mathbf{n}}) 
    \end{align*}
    这里我们将$\delta$作为形式函数来使用，但读者容易对照这个证明，借助线积分给出严格的证明。
\end{proof}

\begin{proposition}[均匀性]
    $\mathcal{R} f(\rho,\phi)$可以延拓为$(\rho,\phi)\in\mathbb{R}^2$的函数，并且
    \[\mathcal{R} f(-\rho,\phi)=\mathcal{R} f(\rho,\phi+\pi)\]
\end{proposition}
\begin{proof}
    形式上，$\delta$是偶分布，因此
    \begin{align*}
        \mathcal{R} f(-\rho,\phi)&=\int_{\mathbb{R}^2}f(x,y)\delta(-\rho - (x\cos\phi+y\sin\phi))\,dx\,dy\\
        &=\int_{\mathbb{R}^2}f(x,y)\delta(\rho + x\cos(\phi+\pi)+y\sin(\phi+\pi))\,dx\,dy\\
        &=\mathcal{R} f(\rho,\phi+\pi)
    \end{align*}
    在拉东变换的分布定义$\left\langle \delta_{L(\rho,\phi)},f\right\rangle$中，$\rho$下，这个命题的严格证明也是容易给出的，不过至此已经能看出拉东变换中使用$\delta$函数的好处。
\end{proof}

\begin{theorem}[投影切片定理]
    将$\mathcal{R} f(\rho,\phi)$关于$\rho$做一维傅里叶变换（这是含参变量$\phi$的积分变换），我们记之为$\mathcal{F} _{\rho}[\mathcal{R} f](r,\phi)$，则有：
    \[\mathcal{F} _{\rho}[\mathcal{R} f](r,\phi)=\mathcal{F} f(r\hat{\mathbf{n}})=\mathcal{F} f(r\cos\phi,r\sin\phi)\]
\end{theorem}
\begin{proof}
    带入定义并交换积分次序：
    \begin{align*}
        \mathcal{F} _{\rho}[\mathcal{R} f](r,\phi)&=\int_{-\infty}^{\infty}\mathcal{R} f(\rho,\phi)\mathrm{e}^{-2\pi \mathrm{i} r \rho}\,d\rho\\
        &=\int_{-\infty}^{\infty}\left(\int_{\mathbb{R}^2}f(x,y)\delta(x\cos\phi+y\sin\phi-\rho)\,dx\,dy\right)\mathrm{e}^{-2\pi \mathrm{i} r \rho}\,d\rho\\
        &=\int_{\mathbb{R}^2}f(x,y)\left(\int_{-\infty}^{\infty}\delta(x\cos\phi+y\sin\phi-\rho)\mathrm{e}^{-2\pi \mathrm{i} r \rho}\,d\rho\right)\,dx\,dy\\
        &=\int_{\mathbb{R}^2}f(x,y)\mathrm{e}^{-2\pi \mathrm{i} r (x\cos\phi+y\sin\phi)}\,dx\,dy\\
        &=\mathcal{F} f(r\cos\phi,r\sin\phi)
    \end{align*}
    要将这个证明严格化，需要限定$f\in\mathcal{S} (\mathbb{R}^2)$，其中$\mathcal{S} (\mathbb{R}^2)$为二维施瓦兹函数空间，将$\delta$用类似\ref{sub:approach}的方法进行逼近，这个过程比较复杂，本书不进行讨论。
\end{proof}

根据投影切片定理，如果我们能够对$[0,\pi)$区间内的每个角度，对足够宽范围内的$\rho$进行测量，就可以得到$\mathcal{R} f(\rho,\phi)$，于是
\[f(x,y)=\mathcal{F} ^{-1}[\mathcal{F} f](x,y)=\mathcal{F} ^{-1}[\mathcal{F} _{\rho}[\mathcal{R} f]](x,y)\]
在充分大的采样频率下，我们可以还原出清晰度足够的图像。这就是x射线成像的基本原理。

\subsection{汉克尔变换}\label{sub:hankel}

如果说多维傅里叶变换是一维傅里叶变换在多维直角坐标系下的版本，那么\textbf{汉克尔变换}(Hankel transform)则是从一维傅里叶变换的“圆形版本”，它适用于具有\textbf{辐射对称}(radial symmetry)或称\textbf{循环对称}(circular symmetry)特性的函数。汉克尔变换需要用到\ref{sub:bessel}中介绍的第一类贝塞尔函数。

\begin{definition}[辐射对称性]
    $f:\mathbb{R}^n\to\mathbb{C}$具有\textbf{辐射对称性}(radial symmetry)，如果
    $$\forall\mathbf{x}\in\mathbb{R}^n,f(\mathbf{x})=f(\|\mathbf{x}\|)$$
    即$f$的极坐标表示仅仅依赖于径向变量$r=\|\mathbf{x}\|$，而与角度无关。
\end{definition}

\begin{definition}[汉克尔变换]
    设$f:[0,\infty)\to\mathbb{C}$，$n\in\mathbb{Z}^+$，则$f$的\textbf{n阶汉克尔变换}(n-th order Hankel transform)定义为：
    \[\mathcal{H} _n f(k)=\int_0^{\infty}f(r)J_n(kr)r\,dr\]
    其中$J_n$为第一类贝塞尔函数。
\end{definition}

本小节主要介绍0阶汉克尔变换，因为它与二维傅里叶变换有着密切的联系，简介起见，本小节使用角频率形式的二维傅里叶变换。

\begin{lemma}
    设$f:\mathbb{R}^2\to\mathbb{C}$具有辐射对称性，则其傅里叶变换$\mathcal{F} f$也具有辐射对称性。
\end{lemma}
\begin{proof}
    为了区别空间域和频率域，我们将空间域的极坐标记为$(r,\theta)$，频率域的极坐标记为$(\rho,\phi)$，于是$f(\mathbf{x})=f(r)$。极坐标下，$\mathbf{x}=(r\cos\theta,r\sin\theta)$，$\boldsymbol{\omega}=(\rho\cos\phi,\rho\sin\phi)$，因此
    $$\boldsymbol{\omega}\cdot\mathbf{x}=\rho r(\cos\theta\cos\phi+\sin\theta\sin\phi)=\rho r \cos(\phi-\theta)$$
    带入二维傅里叶变换的定义，有：
    \begin{align*}
        \mathcal{F} f(\boldsymbol{\omega})&=\int_{\mathbb{R}^2}f(\mathbf{x})\mathrm{e}^{-\mathrm{i}\boldsymbol{\omega}\cdot\mathbf{x}}\,d\mathbf{x}\\
        &=\int_0^{\infty}\int_0^{2\pi}f(r)\mathrm{e}^{-\mathrm{i} \rho r \cos(\phi-\theta)}r\,d\theta\,dr\\
        &=\int_0^{\infty}f(r)r\,dr\int_0^{2\pi}\mathrm{e}^{-\mathrm{i} \rho r \cos(\phi-\theta)}\,d\theta\\
        &=\int_0^{\infty}f(r)r\,dr\int_0^{2\pi}\mathrm{e}^{-\mathrm{i} \rho r \cos\theta}\,d\theta&(\theta\to \theta+\phi)\\
    \end{align*}
    注意到最后一行的表达式与$\phi$无关，因此$\mathcal{F} f(\boldsymbol{\omega})$仅依赖于$\rho$，即具有辐射对称性。
\end{proof}

\begin{theorem}
    设$f:\mathbb{R}^2\to\mathbb{C}$具有辐射对称性，记$f(\mathbf{x})=f(r)$，其中$r=\|\mathbf{x}\|$，则$f$的二维傅里叶变换$\mathcal{F} f(\boldsymbol{\omega})$也具有辐射对称性，记$\mathcal{F} f(\boldsymbol{\omega})=F(\rho)$，其中$\rho=\|\boldsymbol{\omega}\|$，并且有：
    \[F(\rho)=2\pi \int_0^{\infty}f(r)J_0(\rho r)r\,dr=2\pi\mathcal{H}_0 f(\rho)\]
\end{theorem}
\begin{proof}
    根据\ref{sub:bessel}节中的结论，\begin{align*}
        \int_{0}^{2\pi}\mathrm{e}^{-\mathrm{i} \rho r \cos\theta}\,d\theta&=2\pi J_0(\rho r)
    \end{align*}
    带入前一引理的证明中推导的傅里叶变换的表达式，得到：
    \begin{align*}
        \mathcal{F} f(\boldsymbol{\omega})&=\int_0^{\infty}f(r)r\,dr\int_0^{2\pi}\mathrm{e}^{-\mathrm{i} \rho r \cos\theta}\,d\theta\\
        &=2\pi \int_0^{\infty}f(r)J_0(\rho r)r\,dr\\
        &=2\pi \mathcal{H} _0 f(\rho)
    \end{align*}
\end{proof}

\begin{example}[圆形特征函数的傅里叶变换]
    与上一小节相同，记$D=\{(x,y)|x^2+y^2\leq r^2\}$，则它的特征函数$\chi_D(x,y)$的二维傅里叶变换为：
    \begin{align*}
        \mathcal{F} \chi_D(\boldsymbol{\omega})&=2\pi\mathcal{H}_0 \chi_D(\rho)\\
        &=2\pi \int_0^{\infty}\chi_D(r)J_0(\rho r)r\,dr\\
        &=2\pi \int_0^{r}J_0(\rho r)r\,dr
    \end{align*}
    其中$\rho=\|\boldsymbol{\omega}\|$。根据\ref{sub:bessel}节中的结论，有：
    \begin{align*}
        F\chi_D(\boldsymbol{\omega})&=2\pi \left[\frac{r J_1(\rho r)}{\rho}\right]\\
        &=\frac{2\pi r}{\rho}J_1(\rho r)
    \end{align*}
\end{example}

遗憾的是，更高维度下的辐射对称函数的傅里叶变换要复杂得多。事实上，我们有：
\begin{claim}
    设$f:\mathbb{R}^n\to\mathbb{C}$具有辐射对称性，记$f(\mathbf{x})=f(r)$，其中$r=\|\mathbf{x}\|$，则$f$的n维傅里叶变换$\mathcal{F} f(\boldsymbol{\omega})$也具有辐射对称性，记$\mathcal{F} f(\boldsymbol{\omega})=F(\rho)$，其中$\rho=\|\boldsymbol{\omega}\|$，引入常数$\nu=\dfrac{n}{2}-1$，则有：
    \[F(\rho)=(2\pi)^{n/2}\rho^{-\nu}\int_0^{\infty}f(r)J_{\nu}(\rho r)r^{\nu+1}\,dr=(2\pi)^{n/2}\rho^{-\nu}\mathcal{H}_\nu f(\rho)\]
\end{claim}

这意味着，在高维空间中，辐射对称函数的傅里叶变换可以用半整数阶汉克尔变换表示，仅仅使用第一类贝塞尔函数是无法处理奇数维空间的情况的。
